%----------------------------------------------------------------------------------------
% MATRIXMODELLEN. Alle opgaven zijn eigen werk.
%
% Opbouw: Start, leerstof (theorie, werkwijze of uitleg, voorbeeld, korte oefening), en op een
% nieuwe pagina de Opdrachten (door elkaar, stijgende moeilijkheid), Koppel en Test jezelf.
% Labels van dit hoofdstuk beginnen met h4-.
%----------------------------------------------------------------------------------------

\mychapter{4}{Matrixmodellen}

\lkr{Leerplan B+S'': LPD 35 (matrixmodellen gebruiken om evoluties te beschrijven), met de contexten uit de wenken: verbindingsmatrices van een netwerk, overgangsmatrices (Markovketens, migratie) en lesliematrices. Het leerplan laat ICT toe voor hoge machten en vraagt de evolutie op lange termijn te bekijken. Het evenwicht berekenen we exact: bij twee toestanden met één vergelijking, bij drie met de methode van Gauss-Jordan uit hoofdstuk 2. Een hoge macht op het rekentoestel dient als controle. Afspraak in heel het hoofdstuk: de kolom is \textit{van}, de rij is \textit{naar}.}

In de vorige hoofdstukken leerde je rekenen met matrices en stelsels oplossen. Nu gebruiken we matrices om iets te voorspellen. Hoeveel leerlingen eten morgen warm in de refter? Hoeveel fietsen staan er volgende week aan het station? Groeit een populatie uilen, of sterft ze uit?

In al die situaties verandert een toestand stap voor stap, telkens volgens dezelfde regels. Die regels zetten we in een matrix. De toestand na één stap is dan een product: matrix maal kolom. En de toestand na tien stappen vind je met de tiende macht van die matrix. Eerst bekijken we een eenvoudiger vraag: hoeveel wegen er zijn in een netwerk.

\begin{oefeningen}[Start]
\oef{h4-start-1}

\oef{h4-verken}
\end{oefeningen}

\lkr{Laat oefening \ref{oef:h4-verken} maken vóór de uitleg. De leerlingen rekenen met percentages en zien dat het totaal gelijk blijft. In vraag c schrijven ze de berekening zelf als product: zo wordt duidelijk waarom de kolommen bij \textit{van} horen. Vraag d mag een gok blijven; het antwoord volgt in sectie 2.}

%========================================================================================
\section{Grafen en verbindingsmatrices}

\subsection{Een netwerk als matrix}

Een kleine luchtvaartmaatschappij vliegt tussen vier steden: Brussel, Lissabon, Rome en Wenen. Niet elke stad heeft een rechtstreekse vlucht naar elke andere. De figuur toont welke vluchten er zijn: een pijl van Brussel naar Lissabon betekent dat er een vlucht is in die richting.

Zo'n figuur heet een graaf. Ze toont alleen wat er verbonden is, niet hoe ver of in welke richting de steden echt liggen. Om met zo'n netwerk te rekenen, zetten we het in een matrix.

\begin{theorie}
Een \textbf{\textcolor{accent}{graaf}} bestaat uit \textbf{\textcolor{accent}{knopen}} en \textbf{\textcolor{accent}{bogen}} tussen die knopen. Heeft elke boog een richting, dan spreek je van een \textbf{\textcolor{accent}{gerichte graaf}}. Een boog van een knoop naar zichzelf is een \textbf{\textcolor{accent}{lus}}.

\vspace{1mm}
In de \textbf{\textcolor{accent}{verbindingsmatrix}} $V$ van een graaf is $\formule{v_{ij}}$ het aantal bogen \emph{van} knoop $j$ \emph{naar} knoop $i$. De kolommen horen dus bij het vertrek, de rijen bij de aankomst.
\end{theorie}

Waarom staat het vertrek in de kolommen? Bij de modellen in de volgende secties vermenigvuldigen we een matrix met een kolom, en dan moet \textit{van} in de kolommen staan. Die afspraak houden we in heel dit hoofdstuk aan. Let er goed op, want in andere boeken staat het soms omgekeerd.

\voorbeeld{}{Stel de verbindingsmatrix op van het vluchtennet in de figuur.\label{vb:h4-graaf}}{%
\uitleg{We zetten de steden in dezelfde volgorde boven de kolommen en naast de rijen. Dan overlopen we de knopen één voor één: elke pijl die vertrekt, geeft een 1 in de kolom van die knoop, in de rij van de knoop waar de pijl aankomt.}
\naast{\begin{stappenvrij}{4.4cm}{5.2cm}
V = \begin{bmatrix} 0 & 1 & 0 & 1 \\ 1 & 0 & 0 & 0 \\ 1 & 1 & 0 & 1 \\ 0 & 0 & 1 & 0 \end{bmatrix} & De volgorde is $B$, $L$, $R$, $W$, voor de rijen en de kolommen. Uit Brussel vertrekken vluchten naar Lissabon en Rome: in kolom 1 staat een 1 in de rij van $L$ en die van $R$. Uit Lissabon naar Brussel en Rome, uit Rome alleen naar Wenen, uit Wenen naar Brussel en Rome.
\end{stappenvrij}}{\begin{tikzpicture}[font=\small]
  \node[knoop] (b) at (0,2.4) {B};
  \node[knoop] (l) at (0,0) {L};
  \node[knoop] (r) at (2.6,0) {R};
  \node[knoop] (w) at (2.6,2.4) {W};
  \draw[boog] (b) to[bend right=18] (l);
  \draw[boog] (l) to[bend right=18] (b);
  \draw[boog] (b) -- (r);
  \draw[boog] (l) -- (r);
  \draw[boog] (r) to[bend right=18] (w);
  \draw[boog] (w) to[bend right=18] (r);
  \draw[boog] (w) -- (b);
\end{tikzpicture}}}

Een rij van $V$ zegt waar je vandaan kunt komen, een kolom waar je naartoe kunt. Zo heeft kolom 3 maar één 1: uit Rome vertrekt maar één vlucht. Op de hoofddiagonaal staan nullen, want geen enkele vlucht landt in de stad waar ze vertrok.

Is een graaf niet gericht, dan kan je elke boog in twee richtingen volgen. Dan is $v_{ij} = v_{ji}$: de verbindingsmatrix is symmetrisch.

\nujij{Een busnet verbindt drie dorpen: $A$ met $B$, en $B$ met $C$, telkens in beide richtingen. Stel de verbindingsmatrix op.}{$\km{0 & 1 & 0 \\ 1 & 0 & 1 \\ 0 & 1 & 0}$: symmetrisch, want elke lijn telt in beide richtingen.}

\begin{oefeningen}
\oef{h4-graaf}
\end{oefeningen}

\subsection{Wegen in een graaf}

Kun je van Rome naar Lissabon vliegen? Niet rechtstreeks, maar wel met een overstap in Wenen en dan nog een in Brussel. Een opeenvolging van bogen heet een \textbf{weg}, en het aantal bogen is de \textbf{lengte} van de weg. Een vlucht met één overstap is een weg van lengte 2.

Hoeveel wegen van lengte 2 lopen er van $j$ naar $i$? Elke weg gaat via een tussenknoop $k$: eerst van $j$ naar $k$, dan van $k$ naar $i$. Dat is precies wat een element van $V^2$ telt.

\begin{theorie}
\begin{stelling}{Aantal wegen}
Is $V$ de verbindingsmatrix van een graaf, dan is het element in rij $i$ en kolom $j$ van $\formule{V^n}$ het aantal wegen van lengte $n$ van knoop $j$ naar knoop $i$.
\end{stelling}

\begin{bewijs}
Neem eerst $n = 2$. Een weg van lengte 2 van $j$ naar $i$ gaat via een knoop $k$. Voor de eerste boog zijn er $v_{kj}$ mogelijkheden, voor de tweede $v_{ik}$, samen $v_{ik} \cdot v_{kj}$ wegen via $k$. Tel op over alle tussenknopen:
\bewijsregel{v_{i1} v_{1j} + v_{i2} v_{2j} + \cdots + v_{im} v_{mj} = (V^2)_{ij} .}
Een weg van lengte 3 is een weg van lengte 2 gevolgd door één boog. Met dezelfde redenering, nu met $V^2$ en $V$, telt $V^3 = V \cdot V^2$ die wegen, en zo verder voor elke $n$.
\end{bewijs}
\end{theorie}

\voorbeeld{}{Bereken $V^2$ voor het vluchtennet. Hoeveel verbindingen met precies één overstap zijn er van Wenen naar Rome? Tussen welke steden is er geen verbinding met precies één overstap?\label{vb:h4-wegen}}{%
\uitleg{Met één overstap neem je precies twee vluchten: dat is een weg van lengte 2, dus we hebben $V^2$ nodig. Daarna lezen we het element af in de kolom van het vertrek en de rij van de aankomst.}
\begin{stappen}
V^2 = \begin{bmatrix} 1 & 0 & 1 & 0 \\ 0 & 1 & 0 & 1 \\ 1 & 1 & 1 & 1 \\ 1 & 1 & 0 & 1 \end{bmatrix} & Elk element is een rij van $V$ maal een kolom van $V$. Rij $R$ maal kolom $W$ bijvoorbeeld: $1 \cdot 1 + 1 \cdot 0 + 0 \cdot 1 + 1 \cdot 0 = 1$. \\[7ex]
(V^2)_{RW} = 1 & Vertrek in Wenen is kolom 4, aankomst in Rome rij 3. Er is één zo'n verbinding: Wenen, Brussel, Rome. De rechtstreekse vlucht van Wenen naar Rome telt hier niet mee: die heeft lengte 1. \\[1ex]
\text{de nullen} & Met precies één overstap kom je niet van Lissabon of Wenen naar Brussel, niet van Brussel of Rome naar Lissabon, en niet van Rome naar Wenen.
\end{stappen}}

Een 1 op de hoofddiagonaal van $V^2$ is een heen-en-terugreis: in Brussel vertrekken en er met één overstap terugkomen, via Lissabon. Wil je weten hoeveel verbindingen er zijn met hoogstens één overstap, dan tel je $V$ en $V^2$ op.

\begin{werkwijze}
Stel de verbindingsmatrix op, met het vertrek in de kolommen. Voor wegen van lengte $n$ bereken je $V^n$, met de hand of met je rekentoestel. Lees het element af in de kolom van het vertrek en de rij van de aankomst.
\end{werkwijze}

\nujij{Hoeveel verbindingen met precies twee overstappen zijn er van Lissabon naar Brussel? Gebruik $V^3 = V \cdot V^2$.}{Rij $B$ van $V$ maal kolom $L$ van $V^2$: $0 \cdot 0 + 1 \cdot 1 + 0 \cdot 1 + 1 \cdot 1 = 2$. Dat zijn L, B, L, B en L, R, W, B.}

\begin{oefeningen}
\oef{h4-wegen}
\end{oefeningen}

%========================================================================================
\section{Overgangsmatrices}

\subsection{Van toestand naar toestand}

In de refter van een school kunnen de 800 leerlingen elke middag kiezen: een warme maaltijd ($W$) of boterhammen ($B$). Wie vandaag warm eet, doet dat morgen met een kans van 60~\% opnieuw, en neemt met een kans van 40~\% boterhammen. Van wie boterhammen eet, eet morgen 20~\% warm en blijft 80~\% bij de boterhammen.

Die percentages tonen we in een graaf, met bij elke boog de kans. Zo'n graaf heet een \textbf{overgangsdiagram}. De kansen zetten we, met dezelfde afspraak als bij de vluchten, in een matrix.

\figuren{\begin{tikzpicture}[font=\small]
  \node[knoop] (w) at (0,0) {W};
  \node[knoop] (b) at (3.4,0) {B};
  \draw[boog] (w) to[bend left=22] node[kans,above] {0,4} (b);
  \draw[boog] (b) to[bend left=22] node[kans,below] {0,2} (w);
  \draw[boog] (w) to[lus=180] node[kans,left] {0,6} (w);
  \draw[boog] (b) to[lus=0] node[kans,right] {0,8} (b);
\end{tikzpicture}\qquad\qquad
\raisebox{5mm}{\small\begin{tabular}{r|cc} & van $W$ & van $B$ \\ \hline naar $W$ & 0,6 & 0,2 \\ naar $B$ & 0,4 & 0,8 \end{tabular}\quad $\longrightarrow$\quad $M = \begin{bmatrix} 0{,}6 & 0{,}2 \\ 0{,}4 & 0{,}8 \end{bmatrix}$}}

\begin{theorie}
Een systeem dat stap voor stap van de ene toestand naar de andere gaat, met vaste kansen, heet een \textbf{\textcolor{accent}{Markovketen}}. In de \textbf{\textcolor{accent}{overgangsmatrix}} $M$ is $\formule{m_{ij}}$ de kans om \emph{van} toestand $j$ \emph{naar} toestand $i$ te gaan. Alle elementen liggen tussen 0 en 1, en de som van elke kolom is $\formule{1}$.

\vspace{1mm}
De \textbf{\textcolor{accent}{toestandsvector}} $X_n$ is de kolom met de aantallen (of de aandelen) in elke toestand na $n$ stappen.

\begin{stelling}{Matrixmodel}
\qquad $\formule{X_{n+1} = M X_n}$ \qquad en dus \qquad $\formule{X_n = M^n X_0}$
\end{stelling}

\begin{bewijs}
Na één stap is het aantal in toestand $i$ de som over alle toestanden $j$ van (kans van $j$ naar $i$) maal (aantal in $j$), en dat is rij $i$ van $M$ maal de kolom $X_n$. Herhaal je dat, dan is $X_2 = M X_1 = M (M X_0) = M^2 X_0$, en zo verder.
\end{bewijs}
\end{theorie}

De kolomsom is 1, omdat iedereen die vertrekt uit een toestand ergens terechtkomt. Daarom blijft ook het totaal gelijk: er komen geen leerlingen bij en er gaan er geen weg. Zo'n systeem heet \textbf{gesloten}.

\voorbeeld{}{Op maandag eten 300 leerlingen warm en 500 nemen boterhammen. Hoeveel leerlingen eten warm op dinsdag en op woensdag?\label{vb:h4-markov}}{%
\uitleg{We zetten de aantallen van maandag in een kolom, in dezelfde volgorde als de toestanden in $M$: eerst $W$, dan $B$. Elke dag vermenigvuldigen we met $M$.}
\begin{stappenbreed}
X_0 = \begin{bmatrix} 300 \\ 500 \end{bmatrix} & De toestand op maandag. \\[3.5ex]
X_1 = M X_0 = \begin{bmatrix} 0{,}6 \cdot 300 + 0{,}2 \cdot 500 \\ 0{,}4 \cdot 300 + 0{,}8 \cdot 500 \end{bmatrix} = \begin{bmatrix} 280 \\ 520 \end{bmatrix} & Op dinsdag eten warm: 60~\% van wie maandag warm at, plus 20~\% van wie boterhammen at. \\[3.5ex]
X_2 = M X_1 = \begin{bmatrix} 0{,}6 \cdot 280 + 0{,}2 \cdot 520 \\ 0{,}4 \cdot 280 + 0{,}8 \cdot 520 \end{bmatrix} = \begin{bmatrix} 272 \\ 528 \end{bmatrix} & Op woensdag: dezelfde regels, nu op de toestand van dinsdag. Het totaal blijft 800.
\end{stappenbreed}}

\begin{werkwijze}
Teken het overgangsdiagram en zet de kansen in de overgangsmatrix: van in de kolommen, naar in de rijen. Controleer dat elke kolomsom 1 is. Zet de begintoestand in een kolom, in dezelfde volgorde, en vermenigvuldig: $X_1 = M X_0$, $X_2 = M X_1$, of meteen $X_n = M^n X_0$.
\end{werkwijze}

\nujij{Hoeveel leerlingen eten warm op donderdag?}{$X_3 = M X_2$: $0{,}6 \cdot 272 + 0{,}2 \cdot 528 = 268{,}8$, dus ongeveer 269.}

\begin{oefeningen}
\oef{h4-overgang}
\end{oefeningen}

Een model met drie toestanden werkt op dezelfde manier, maar het rekenwerk groeit. Dan laat je het rekentoestel de machten berekenen.

\voorbeeld{}{Een stad heeft deelfietsen aan drie stations: het Station ($S$), de Markt ($M$) en de Campus ($C$). Elke dag komen de fietsen aan een station terecht volgens het diagram. Op maandagochtend staan er 120 fietsen aan het Station, 80 aan de Markt en 100 aan de Campus. Hoeveel fietsen staan er een week later aan elk station?\label{vb:h4-fietsen}}{%
\uitleg{Eerst lezen we de overgangsmatrix af uit het diagram, kolom per vertrekstation. Een week later is zeven stappen verder, dus we hebben $F^7$ nodig, en dat laten we aan het rekentoestel over.}
\figuren{\begin{tikzpicture}[font=\small]
  \node[knoop] (s) at (0,0) {S};
  \node[knoop] (m) at (3.4,0) {M};
  \node[knoop] (c) at (1.7,-2.9) {C};
  \draw[boog] (s) to[bend left=18] node[kans,above] {0,3} (m);
  \draw[boog] (m) to[bend left=18] node[kans,below] {0,2} (s);
  \draw[boog] (s) to[bend left=18] node[kans,pos=.45,right=1pt] {0,2} (c);
  \draw[boog] (c) to[bend left=18] node[kans,pos=.55,left=1pt] {0,3} (s);
  \draw[boog] (m) to[bend left=18] node[kans,pos=.45,right=1pt] {0,2} (c);
  \draw[boog] (c) to[bend left=18] node[kans,pos=.55,left=1pt] {0,2} (m);
  \draw[boog] (s) to[lus=180] node[kans,left] {0,5} (s);
  \draw[boog] (m) to[lus=0] node[kans,right] {0,6} (m);
  \draw[boog] (c) to[lus=270] node[kans,below] {0,5} (c);
\end{tikzpicture}}
\begin{stappen}
F = \begin{bmatrix} 0{,}5 & 0{,}2 & 0{,}3 \\ 0{,}3 & 0{,}6 & 0{,}2 \\ 0{,}2 & 0{,}2 & 0{,}5 \end{bmatrix} & Kolom 1: van het Station blijft 50~\%, 30~\% gaat naar de Markt en 20~\% naar de Campus. Elke kolomsom is 1. \\[5ex]
X_7 = F^7 \begin{bmatrix} 120 \\ 80 \\ 100 \end{bmatrix} \approx \begin{bmatrix} 98{,}0 \\ 116{,}3 \\ 85{,}7 \end{bmatrix} & Met het rekentoestel. Er staan dus ongeveer 98 fietsen aan het Station, 116 aan de Markt en 86 aan de Campus.
\end{stappen}}

\begin{oefeningen}
\oef{h4-drie}
\end{oefeningen}

\subsection{Op lange termijn}

Rekenen we in de refter verder, dan verandert het aantal warme maaltijden steeds minder: 280, 272, 268,8, \ldots{} Het lijkt naar een vaste waarde te gaan. En dat gebeurt ook als je op maandag met heel andere aantallen begint. De figuur toont het aandeel leerlingen dat warm eet, voor twee verschillende begintoestanden: na een week of twee zijn ze niet meer van elkaar te onderscheiden.

\figuren{\begin{tikzpicture}[x=.62cm,y=4.2cm,font=\scriptsize]
  \draw[black!14,very thin] (0,0) grid[xstep=1,ystep=.1] (10,1);
  \draw[->,grijs] (0,0) -- (10.6,0) node[right,text=black] {dag $n$};
  \draw[->,grijs] (0,0) -- (0,1.08) node[above,text=black] {aandeel warm};
  \foreach \t in {0,2,...,10} \node[below] at (\t,0) {\t};
  \foreach \t/\l in {.2/{0,2},.4/{0,4},.6/{0,6},.8/{0,8},1/1} \node[left] at (0,\t) {\l};
  \draw[accent,dashed] (0,1/3) -- (10,1/3) node[right,text=accent] {$\tfrac13$};
  \draw[accent,thick] plot[mark=*,mark size=1.3pt] coordinates {(0,.375) (1,.35) (2,.34) (3,.336) (4,.3344) (5,.3338) (6,.3335) (7,.3334) (8,.3333) (9,.3333) (10,.3333)};
  \draw[fel!80!black,thick] plot[mark=square*,mark size=1.2pt] coordinates {(0,.875) (1,.55) (2,.42) (3,.368) (4,.3472) (5,.3389) (6,.3356) (7,.3342) (8,.3337) (9,.3335) (10,.3334)};
  \node[right,text=fel!80!black] at (.3,.875) {start: 700 warm};
  \node[above,text=accent] at (2.6,.43) {start: 300 warm};
\end{tikzpicture}}

Op lange termijn eet dus een derde van de leerlingen warm, welke dag je ook begint. Elke dag wisselen er leerlingen van keuze, maar evenveel in elke richting: 40~\% van de warme eters is evenveel als 20~\% van de boterhameters, als die tweede groep dubbel zo groot is. Zo'n toestand verandert niet meer.

\begin{theorie}
Een \textbf{\textcolor{accent}{evenwicht}} van een overgangsmatrix $M$ is een toestand $X$ die niet meer verandert: \quad $\formule{MX = X}$
\end{theorie}

Hoe vind je het evenwicht? Bij twee toestanden lukt het met één vergelijking. We werken met aandelen: $x$ eet warm en $1 - x$ neemt boterhammen.

\voorbeeld{}{Bereken het evenwicht in de refter.\label{vb:h4-evenwicht}}{%
\uitleg{We zoeken de aandelen die gelijk blijven na één stap. Het aandeel warm na een dag moet dus opnieuw $x$ zijn. Dat geeft een eerstegraadsvergelijking in $x$.}
\begin{stappen}
\begin{bmatrix} 0{,}6 & 0{,}2 \\ 0{,}4 & 0{,}8 \end{bmatrix} \begin{bmatrix} x \\ 1 - x \end{bmatrix} = \begin{bmatrix} x \\ 1 - x \end{bmatrix} & De voorwaarde $MX = X$, met $x + (1 - x) = 1$. \\[3.5ex]
0{,}6x + 0{,}2(1 - x) = x & De eerste rij is genoeg: als die klopt, klopt de tweede ook, want het totaal blijft 1. \\[1ex]
0{,}2 = 0{,}6x \ \Rightarrow \ x = \tfrac13 & Werk de haakjes uit: $0{,}4x + 0{,}2 = x$. \\[1.5ex]
X = \begin{bmatrix} \tfrac13 \\[.5ex] \tfrac23 \end{bmatrix} & Op lange termijn eet een derde warm, ongeveer 267 leerlingen, en twee derde boterhammen, ongeveer 533.
\end{stappen}}

Bij drie toestanden geeft $MX = X$ een stelsel met drie onbekenden, dat je oplost met de methode van Gauss-Jordan uit hoofdstuk 2. Schrijf het als $(M - I)X = O$. Dat stelsel is onbepaald: is $X$ een evenwicht, dan ook $2X$. De voorwaarde dat de aandelen samen 1 zijn, legt de oplossing vast.

\voorbeeld{}{Bepaal het evenwicht van de deelfietsen uit voorbeeld \ref{vb:h4-fietsen}.\label{vb:h4-lang}}{%
\uitleg{We lossen $(F - I)X = O$ op met Gauss-Jordan. Om met gehele getallen te rekenen, vermenigvuldigen we elke vergelijking eerst met 10. Daarna kiezen we de parameter zo dat de aandelen samen 1 zijn.}
\begin{stappenbreed}
\left[\begin{array}{ccc|c} -5 & 2 & 3 & 0 \\ 3 & -4 & 2 & 0 \\ 2 & 2 & -5 & 0 \end{array}\right] & $10(F - I)X = O$: de uitgebreide matrix. \\[5ex]
\left[\begin{array}{ccc|c} 1 & 0 & -\frac87 & 0 \\[.3ex] 0 & 1 & -\frac{19}{14} & 0 \\[.3ex] 0 & 0 & 0 & 0 \end{array}\right] & Na rijherleiden, of met \texttt{rref} op het rekentoestel. Eén nulrij: onbepaald. \\[6ex]
x = \tfrac{16}{14}z, \quad y = \tfrac{19}{14}z & Uit de eerste twee rijen, met $z$ als parameter. \\[1.5ex]
\tfrac{16}{14}z + \tfrac{19}{14}z + \tfrac{14}{14}z = 1 \ \Rightarrow \ z = \tfrac{14}{49} & De aandelen zijn samen 1. \\[1.5ex]
X = \begin{bmatrix} 16/49 \\ 19/49 \\ 14/49 \end{bmatrix} \approx \begin{bmatrix} 0{,}327 \\ 0{,}388 \\ 0{,}286 \end{bmatrix} & Met 300 fietsen: ongeveer 98 aan het Station, 116 aan de Markt en 86 aan de Campus. Een week na de start zat het aantal dus al zo goed als op het evenwicht.
\end{stappenbreed}
Controle met het rekentoestel: alle kolommen van $F^{30}$ zijn op drie decimalen gelijk aan $X$.}

Dat de kolommen van een hoge macht gelijk worden, is geen toeval: $M^{n} X_0$ gaat naar het evenwicht voor elke begintoestand, dus ook voor de begintoestand met alles in toestand 1, en dan is $M^n X_0$ precies de eerste kolom van $M^n$. Je zag dat al in hoofdstuk 1.

\begin{werkwijze}
Bij twee toestanden: los $MX = X$ op met $X = \km{x \\ 1 - x}$, met één rij. Bij meer toestanden: los $(M - I)X = O$ op met Gauss-Jordan, en kies de parameter zo dat de som van de aandelen 1 is. Vermenigvuldig met het totaal voor de aantallen. Controleer met een hoge macht van $M$: elke kolom is dan het evenwicht.
\end{werkwijze}

Niet elke Markovketen komt zo in een evenwicht met alle toestanden. Zit er een toestand in die je nooit meer verlaat, zoals een abonnement dat is opgezegd, dan komt op lange termijn alles daar terecht.

\begin{oefeningen}
\oef{h4-evenwicht}

\oef{h4-evenwicht3}
\end{oefeningen}

%========================================================================================
\section{Lesliematrices}

\subsection{Een populatie in leeftijdsklassen}

Hoe een populatie dieren evolueert, hangt af van haar leeftijd: jonge dieren krijgen nog geen jongen, en oude dieren sterven sneller. Daarom verdelen we de populatie in \textbf{leeftijdsklassen} die even lang duren als één stap van het model. Bij dieren tellen we meestal alleen de vrouwtjes, want het aantal jongen hangt van hen af.

We bekijken een fictieve populatie kerkuilen in drie klassen van één jaar: 0 tot 1 jaar, 1 tot 2 jaar en 2 tot 3 jaar. Van de uilen in klasse 1 overleeft 40~\% het jaar, en schuift door naar klasse 2. Van klasse 2 overleeft 60~\%. Een vrouwtje in klasse 2 krijgt in een jaar gemiddeld 1,5 vrouwelijke jongen, een vrouwtje in klasse 3 gemiddeld 2. Een uil in klasse 3 wordt niet ouder dan 3 jaar.

\naast{\begin{minipage}{9.4cm}
Er zijn nu twee soorten pijlen. De pijlen naar rechts zijn overlevingskansen: een uil schuift een klasse op. De pijlen terug naar klasse 1 zijn de jongen die geboren worden. Ook dat zetten we in een matrix, met de vertrekklasse in de kolommen.

\vspace{1mm}
\qquad $L = \begin{bmatrix} 0 & 1{,}5 & 2 \\ 0{,}4 & 0 & 0 \\ 0 & 0{,}6 & 0 \end{bmatrix}$
\end{minipage}}{\begin{tikzpicture}[font=\small]
  \node[knoop] (k1) at (0,0) {1};
  \node[knoop] (k2) at (2.4,0) {2};
  \node[knoop] (k3) at (4.8,0) {3};
  \draw[boog] (k1) -- node[kans,below=1pt] {0,4} (k2);
  \draw[boog] (k2) -- node[kans,below=1pt] {0,6} (k3);
  \draw[boog,dashed] (k2) to[out=120,in=60] node[kans,above] {1,5} (k1);
  \draw[boog,dashed] (k3) to[out=110,in=80,looseness=.9] node[kans,above] {2} (k1);
\end{tikzpicture}}

\begin{theorie}
Een \textbf{\textcolor{accent}{lesliematrix}} $L$ beschrijft een populatie in leeftijdsklassen die even lang duren als één stap. In de eerste rij staan de \textbf{\textcolor{accent}{vruchtbaarheidscijfers}}: het gemiddeld aantal jongen per dier in elke klasse. Net onder de hoofddiagonaal staan de \textbf{\textcolor{accent}{overlevingskansen}}: de kans dat een dier doorschuift naar de volgende klasse. Alle andere elementen zijn 0.

\vspace{1mm}
\qquad $\formule{L = \begin{bmatrix} v_1 & v_2 & v_3 \\ o_1 & 0 & 0 \\ 0 & o_2 & 0 \end{bmatrix}}$
\end{theorie}

Ook hier geldt $X_{n+1} = L X_n$ en dus $X_n = L^n X_0$: de redenering is dezelfde als bij een overgangsmatrix. Er is wel een groot verschil. De kolommen van een lesliematrix hebben geen som 1, omdat er dieren sterven en geboren worden. Het totaal kan dus groeien of krimpen.

\voorbeeld{}{Bij de start zijn er 100 uilen in klasse 1, 60 in klasse 2 en 40 in klasse 3. Bereken de toestand na één en na twee jaar.\label{vb:h4-leslie}}{%
\uitleg{We vermenigvuldigen de lesliematrix met de kolom van de begintoestand. De eerste rij geeft de jongen van dit jaar, de andere rijen de dieren die een klasse opschuiven.}
\begin{stappen}
X_1 = L \begin{bmatrix} 100 \\ 60 \\ 40 \end{bmatrix} = \begin{bmatrix} 1{,}5 \cdot 60 + 2 \cdot 40 \\ 0{,}4 \cdot 100 \\ 0{,}6 \cdot 60 \end{bmatrix} = \begin{bmatrix} 170 \\ 40 \\ 36 \end{bmatrix} & Klasse 1 bestaat nu uit de pasgeboren jongen: 90 van klasse 2 en 80 van klasse 3. Van de 100 jonge uilen leven er nog 40, die nu in klasse 2 zitten. \\[5ex]
X_2 = L \begin{bmatrix} 170 \\ 40 \\ 36 \end{bmatrix} = \begin{bmatrix} 60 + 72 \\ 68 \\ 24 \end{bmatrix} = \begin{bmatrix} 132 \\ 68 \\ 24 \end{bmatrix} & Hetzelfde met de toestand na één jaar. Het totaal ging van 200 naar 246 en dan naar 224.
\end{stappen}

\figuren{\begin{tikzpicture}[x=.36cm,y=.016cm,font=\scriptsize]
  \foreach \j/\a/\b/\c in {0/100/60/40,1/170/40/36,2/132/68/24,3/150/52.8/40.8}{
    \begin{scope}[shift={(\j*4.6,0)}]
      \fill[accent!70] (0,0) rectangle (1,\a);
      \fill[accent!40] (1.1,0) rectangle (2.1,\b);
      \fill[accent!15] (2.2,0) rectangle (3.2,\c);
      \draw[accent] (0,0) rectangle (1,\a) (1.1,0) rectangle (2.1,\b) (2.2,0) rectangle (3.2,\c);
      \node[below] at (1.6,0) {jaar \j};
    \end{scope}}
  \draw[lijngrijs] (-.4,0) -- (17.8,0);
  \foreach \t in {50,100,150} {\draw[lijngrijs] (-.5,\t) -- (-.3,\t); \node[left] at (-.4,\t) {\t};}
  \node[right] at (18.2,150) {\tikz\fill[accent!70] (0,0) rectangle (2mm,2mm); klasse 1};
  \node[right] at (18.2,110) {\tikz\fill[accent!40] (0,0) rectangle (2mm,2mm); klasse 2};
  \node[right] at (18.2,70) {\tikz\fill[accent!15] (0,0) rectangle (2mm,2mm); klasse 3};
\end{tikzpicture}}}

\begin{werkwijze}
Zet de vruchtbaarheidscijfers in de eerste rij en de overlevingskansen net onder de hoofddiagonaal. Controleer dat een stap even lang duurt als een leeftijdsklasse. Vermenigvuldig dan met de begintoestand: $X_1 = L X_0$, of $X_n = L^n X_0$ met het rekentoestel.
\end{werkwijze}

\nujij{Bereken het aantal uilen in klasse 3 na drie jaar, zonder de hele $X_3$ uit te rekenen.}{Rij 3 van $L$ maal $X_2$: $0{,}6 \cdot 68 = 40{,}8$, dus ongeveer 41 uilen.}

\begin{oefeningen}
\oef{h4-leslie}
\end{oefeningen}

\subsection{Groei op lange termijn}

In de eerste jaren springt het aantal uilen op en neer: 200, 246, 224, 244, \ldots{} Dat komt door de begintoestand, met veel jonge dieren die nog geen jongen krijgen. Rekenen we verder met het rekentoestel, dan verdwijnt dat schokken. De figuur toont het totaal over 25 jaar.

\figuren{\begin{tikzpicture}[x=.44cm,y=.0075cm,font=\scriptsize]
  \draw[black!14,very thin] (0,0) grid[xstep=5,ystep=100] (25,500);
  \draw[->,grijs] (0,0) -- (25.8,0) node[right,text=black] {jaar $n$};
  \draw[->,grijs] (0,0) -- (0,540) node[above,text=black] {aantal uilen};
  \foreach \t in {0,5,...,25} \node[below] at (\t,0) {\t};
  \foreach \t in {100,200,...,500} \node[left] at (0,\t) {\t};
  \draw[accent,thick] plot[mark=*,mark size=1.1pt] coordinates {(0,200) (1,246) (2,224) (3,243.6) (4,252.5) (5,253.7) (6,268.4) (7,273.4) (8,282.8) (9,292.9) (10,300.9) (11,311.5) (12,321.1) (13,331.3) (14,342.2) (15,352.9) (16,364.4) (17,376.0) (18,388.0) (19,400.5) (20,413.3) (21,426.6) (22,440.2) (23,454.3) (24,468.9) (25,483.9)};
\end{tikzpicture}}

Na een tijd neemt het totaal elk jaar met hetzelfde percentage toe: de populatie groeit exponentieel. Ook de verhouding tussen de klassen ligt dan vast.

\voorbeeld{}{Bereken met je rekentoestel $X_{20}$ en $X_{21}$ voor de uilen. Met welke factor groeit de populatie per jaar op lange termijn? Hoe is de populatie dan over de klassen verdeeld?\label{vb:h4-groei}}{%
\uitleg{We vergelijken het totaal in twee opeenvolgende jaren: hun quotiënt is de groeifactor. Voor de verdeling delen we elke klasse door het totaal.}
\begin{stappen}
X_{20} \approx \begin{bmatrix} 256{,}2 \\ 99{,}3 \\ 57{,}7 \end{bmatrix}, \quad X_{21} \approx \begin{bmatrix} 264{,}5 \\ 102{,}5 \\ 59{,}6 \end{bmatrix} & Bereken $L^{20} X_0$ en $L^{21} X_0$. \\[5ex]
\dfrac{426{,}6}{413{,}3} \approx 1{,}032 & De totalen. Per jaar komt er ongeveer 3,2~\% bij. Ook elke klasse apart groeit met die factor: $\tfrac{264{,}5}{256{,}2} \approx 1{,}032$. \\[2.5ex]
\dfrac{256{,}2}{413{,}3} \approx 0{,}62 & In klasse 1 zit 62~\% van de uilen, in klasse 2 ongeveer 24~\% en in klasse 3 ongeveer 14~\%.
\end{stappen}}

Een groeifactor groter dan 1 betekent dat de populatie op lange termijn groeit, een factor kleiner dan 1 dat ze uitsterft. Bij een factor van precies 1 blijft ze stabiel. Soms ontstaat er geen vaste factor, maar herhaalt de populatie zich om de zoveel jaar: dat zie je in de opdrachten.

\begin{werkwijze}
Bereken met het rekentoestel de toestand na veel stappen en na één stap meer. Het quotiënt van de totalen is de groeifactor. Verandert die factor niet meer bij nog meer stappen, dan is dat de groei op lange termijn.
\end{werkwijze}

\begin{oefeningen}
\oef{h4-groei}
\end{oefeningen}

%========================================================================================
\opdrachtenkop

\begin{oefeningen}
\needspace{14\baselineskip}
\oef{h4-o1}

\oef{h4-o4}

\oef{h4-o2}

\needspace{14\baselineskip}
\oef{h4-o20}

\needspace{16\baselineskip}
\oef{h4-o5}

\needspace{14\baselineskip}
\oef{h4-o3}

\oef{h4-o21}

\needspace{12\baselineskip}
\oef{h4-o6}

\oef{h4-o7}

\needspace{14\baselineskip}
\oef{h4-o8}

\needspace{14\baselineskip}
\oef{h4-o22}

\oef{h4-o11}

\oef{h4-o9}

\needspace{14\baselineskip}
\oef{h4-o23}

\oef{h4-o10}

\oef{h4-o14}

\oef{h4-o25}

\oef{h4-o13}

\oef{h4-o24}

\oef{h4-o12}

\oef{h4-o15}

\oef{h4-o16}

\oef{h4-o26}

\oef{h4-o19}

\oef{h4-o17}

\oef{h4-o18}
\end{oefeningen}

%----------------------------------------------------------------------------------------
% KOPPEL: verbanden met andere leerstof
%----------------------------------------------------------------------------------------
\needspace{20\baselineskip}
\opdrachtsectie{Koppel}

\lkr{Kansrekenen, Rijen en Exponentiële en logaritmische functies zijn cursussen van hetzelfde pakket. Het pakket legt hun volgorde niet vast: kijk na wat de leerlingen al zagen. Voor de laatste opdracht hebben ze logaritmen nodig, of het rekentoestel.}

\begin{oefeningen}
\oef{h4-koppel-1}

\oef{h4-koppel-2}

\oef{h4-koppel-3}
\end{oefeningen}

%----------------------------------------------------------------------------------------
\opdrachtsectie{Test jezelf}

\begin{oefeningen}
\oef{h4-t1}

\oef{h4-t2}

\oef{h4-t3}

\oef{h4-t4}

\oef{h4-t5}

\oef{h4-t6}

\oef{h4-t7}
\end{oefeningen}

\herhaalschema{%
  \fv{oef:h4-t1}{vb:h4-graaf}{oef:h4-graaf}%
  \fv{oef:h4-t2}{vb:h4-wegen}{oef:h4-wegen}%
  \fv{oef:h4-t3}{vb:h4-markov}{oef:h4-overgang}%
  \fv{oef:h4-t4}{vb:h4-evenwicht}{oef:h4-evenwicht}%
  \fv{oef:h4-t5}{vb:h4-leslie}{oef:h4-leslie}%
  \fv{oef:h4-t6}{vb:h4-groei}{oef:h4-groei}%
  \fv{oef:h4-t7}{vb:h4-lang}{oef:h4-evenwicht3}}

\ifzelfstudie
\section*{Maak je eigen samenvatting}
Sluit dit boek en schrijf in je schrift: wat een graaf en een verbindingsmatrix zijn, wat de machten van die matrix tellen, wat een overgangsmatrix is en hoe je er de volgende toestand mee berekent, wat een evenwicht is en hoe je het vindt, en hoe een lesliematrix eruitziet. Schrijf ook op welk verschil er is tussen een overgangsmatrix en een lesliematrix. Vergelijk daarna met de kaders.
\fi
